Il progetto integra fonti operative e cataloghi tematici Copernicus. GDACS,
CDSE/STAC e CEMS On-demand Mapping descrivono eventi e allerte, acquisizioni
satellitari e attivazioni cartografiche di emergenza. Marine e Drought ampliano
la scoperta di risorse attraverso un catalogo di metadati a grafo. Ogni fonte
risponde a una domanda diversa: una heatmap di acquisizioni, per esempio, mostra
dove si concentrano le immagini disponibili, mentre una prima stima dell'impatto di un
disastro viene dalle allerte GDACS. Il capitolo descrive quindi contenuto, modalità di
accesso e limiti di ciascuna fonte prima di esaminare l'architettura che le
combina.

\section{GDACS}

\subsection{Scopo e contesto operativo}

GDACS, acronimo di \emph{Global Disaster Alert and Coordination System}, è un
quadro di cooperazione tra le Nazioni Unite, la Commissione europea e gli attori
internazionali della gestione delle emergenze. Il suo obiettivo principale è
ridurre le carenze informative e di coordinamento nella prima fase successiva a
grandi disastri improvvisi \cite{gdacsOverview}.

La piattaforma raccoglie dati sugli eventi, esegue o riceve stime automatiche
dell'impatto potenziale e distribuisce notifiche. La sua funzione di
\emph{early awareness} mette in relazione caratteristiche fisiche dell'evento,
popolazione esposta, vulnerabilità e capacità di risposta. I modelli variano in
base al tipo di pericolo e producono una valutazione sintetica utile a stabilire
quali eventi richiedano attenzione internazionale \cite{gdacsKnowledge}.

\subsection{Eventi e livelli di allerta}

GDACS pubblica informazioni su diverse categorie di pericolo, tra cui terremoti,
possibili tsunami, cicloni tropicali, alluvioni, eruzioni vulcaniche, siccità e
incendi. I record disponibili contengono normalmente:

\begin{itemize}
    \item identificativo e tipologia dell'evento;
    \item titolo, paese e coordinate;
    \item istante di pubblicazione o aggiornamento;
    \item livello di allerta;
    \item misure specifiche, come magnitudo o severità;
    \item collegamento alla scheda ufficiale.
\end{itemize}

I livelli verde, arancione e rosso esprimono una stima crescente dell'impatto e
della rilevanza internazionale. Oltre alla misura fisica del fenomeno, il livello
considera esposizione, vulnerabilità e capacità di risposta: un forte terremoto
in un'area poco popolata può quindi ricevere un'allerta bassa.

\subsection{Accesso programmatico}

I dati e le stime GDACS sono accessibili tramite feed RSS e altri formati
standard, così da poter essere integrati in portali e applicazioni esterne
\cite{gdacsOverview}. Sono disponibili feed distinti per intervallo temporale e
tipologia di evento; la documentazione indica, ad esempio, feed relativi alle
ultime 24 ore, all'ultima settimana e a specifiche soglie per i terremoti
\cite{gdacsFeeds}.

Nel progetto il feed viene interrogato da un servizio MCP autonomo, che
normalizza il contenuto XML e restituisce oggetti strutturati. L'agente riceve
così un contratto stabile e filtra i record per luogo, periodo, tipo di evento
e, quando richiesto, soglie quantitative.

\subsection{Limiti}

Una notifica GDACS è una stima rapida e può cambiare quando arrivano nuovi
dati; le valutazioni sul campo e i bollettini ufficiali di protezione civile
restano il riferimento definitivo. Alcuni processi sono automatici, altri
includono dati e valutazioni provenienti da istituzioni o fonti esterne. Per
questo il progetto conserva link, timestamp e identificativi: ogni risposta
dell'agente resta verificabile e presenta le stime preliminari come tali.

\section{Il programma Copernicus}

Copernicus è il programma di osservazione della Terra dell'Unione europea.
Mette a disposizione missioni satellitari, dati, infrastrutture di accesso e
servizi tematici. Per questa tesi sono rilevanti percorsi complementari:

\begin{enumerate}
    \item il \emph{Copernicus Data Space Ecosystem} (CDSE), usato per cercare
    metadati relativi alle acquisizioni satellitari;
    \item il \emph{Copernicus Emergency Management Service} (CEMS), usato per
    consultare attivazioni e prodotti cartografici di emergenza;
    \item Copernicus Marine, per individuare prodotti, dataset e variabili
    oceanografiche;
    \item i servizi Drought consultati dall'harvester, per scoprire layer,
    coverage e disponibilità temporali relativi alla siccità.
\end{enumerate}

Il catalogo satellitare descrive osservazioni disponibili; CEMS documenta
invece attivazioni operative e prodotti derivati dall'analisi di immagini e di
altre informazioni.

\section{Copernicus Data Space Ecosystem e STAC}

\subsection{Ricerca delle acquisizioni}

Il CDSE offre API per catalogazione, download, visualizzazione e processamento
dei dati di osservazione della Terra. Fra le API di catalogo è disponibile
STAC, insieme a OData e Sentinel Hub Catalog API \cite{cdseApis}.

STAC, \emph{SpatioTemporal Asset Catalog}, definisce un modello interoperabile
per descrivere collezioni e singoli elementi geospaziali. Un elemento associa
un identificativo a geometria, intervallo temporale, proprietà, collegamenti e
asset \cite{stacSpec}. La ricerca può quindi combinare:

\begin{itemize}
    \item una o più collezioni, ad esempio Sentinel-2 Level-2A;
    \item un intervallo temporale;
    \item un'area espressa come bounding box;
    \item filtri su attributi, quali copertura nuvolosa o piattaforma;
    \item un limite al numero di risultati.
\end{itemize}

Il catalogo CDSE implementa STAC 1.1.0 ed espone un endpoint di ricerca
standardizzato \cite{cdseStac}.

\subsection{Dal metadato alla heatmap}

Il ramo STAC del progetto lavora sui metadati: recupera le acquisizioni che
intersecano luogo e periodo richiesti, senza analizzare il contenuto delle
immagini. Di ogni prodotto si usa il rettangolo che lo racchiude (bounding box),
se ne calcola il centro e si costruisce:

\begin{itemize}
    \item una heatmap che mostra dove si concentrano le acquisizioni recuperate;
    \item i rettangoli delle bounding box, selezionabili;
    \item un pannello con identificativo, collezione, data, piattaforma,
    copertura nuvolosa, asset e collegamenti disponibili.
\end{itemize}

I conteggi mostrati si riferiscono ai prodotti restituiti dalla ricerca: più
prodotti possono derivare dalla stessa ripresa satellitare, ad esempio a
diversi livelli di elaborazione. La visualizzazione descrive quindi la
disponibilità di dati sull'area, come punto di partenza per analisi successive.

\section{Copernicus Emergency Management Service}

CEMS fornisce informazioni geospaziali a supporto della risposta alle
emergenze e della gestione del rischio. È uno dei sei servizi Copernicus e
comprende monitoraggio e allerta precoce, cartografia on demand e informazioni
sull'esposizione \cite{copernicusEMSOverview}.

Nel progetto viene usata la componente \emph{On-demand Mapping}. Le attivazioni
sono identificate da codici come \texttt{EMSR884} o \texttt{EMSN204}. Il
prefisso consente di distinguere due percorsi principali:

\begin{description}
    \item[Rapid Mapping (\texttt{EMSR}).] Produce in tempi rapidi informazioni
    geospaziali per la fase immediatamente successiva a un evento. Una
    attivazione può comprendere più aree di interesse e più prodotti
    \cite{copernicusRapidMapping}.
    \item[Risk and Recovery Mapping (\texttt{EMSN}).] Supporta prevenzione,
    preparazione e ricostruzione mediante analisi geospaziali più articolate,
    basate su immagini ottiche e radar \cite{copernicusRiskRecovery}.
\end{description}

Il relativo server MCP espone operazioni per elencare le attivazioni, ottenere
una sintesi e recuperare i dettagli di una specifica attivazione. I risultati
normalizzati includono codice, nome, categoria, paesi, fase di gestione del
rischio, stato, data, aree di interesse e prodotti disponibili.

\section{Marine e Drought nel catalogo di Discovery}
\label{sec:fonti-discovery}

La richiesta di una risorsa scientifica richiede metadati più articolati del
solo elenco delle acquisizioni. Per esempio, studiare la temperatura
superficiale del Mar Nero implica individuare un dataset, la variabile
disponibile, l'area coperta e la risoluzione temporale. La Discovery usa un
catalogo locale alimentato da tre harvester, documentati nella guida tecnica
del progetto \cite{guidaProgetto2026}.

\subsection{Metadati Marine}

Il collector Marine legge il catalogo completo di Copernicus Marine, comprese
le versioni precedenti dei dataset. Conserva la gerarchia tra prodotto,
dataset, versione, parte e servizio di accesso, oltre a variabili e coordinate.
Una variabile è descritta da nome, nome standard, unità e descrizione; gli assi
da estremi, passo e numero di valori, senza memorizzare i valori completi.

Un secondo collector documenta le operazioni offerte dalla libreria ufficiale
di Copernicus Marine e i relativi parametri. Il catalogo può così indicare,
per ogni risorsa, le modalità con cui sarebbe possibile estrarne una porzione;
l'harvester registra queste possibilità senza scaricare i file
scientifici.

\subsection{Metadati Drought}

Il collector Drought interroga le descrizioni standard dei servizi cartografici
WMS e WCS. Raccoglie servizi, layer e sotto-layer, coverage ed estensioni
temporali. Queste descrizioni consentono di verificare se una risorsa sia
compatibile con tema e periodo richiesti senza scaricare mappe o raster. Il
censimento copre le coverage previste dal catalogo applicativo.

\subsection{Due usi distinti di STAC}

L'harvester CDSE salva root del catalogo, collezioni, campi interrogabili e,
quando previsto, item e riferimenti agli asset. La finestra recente
predefinita per gli item è di due giorni, con limiti configurabili. Degli asset
si conservano i collegamenti, mentre i file restano sul servizio remoto.

Nella Discovery l'agente legge questi metadati da NebulaGraph; nel percorso
operativo il client interroga direttamente CDSE STAC per cercare acquisizioni.
I due percorsi sono indipendenti: i risultati operativi riflettono lo stato
attuale del servizio remoto, mentre il grafo riflette l'ultimo aggiornamento
del catalogo locale.

\section{Confronto e complementarità}

\begin{table}[htbp]
\centering
\small
\begin{tabularx}{\textwidth}{
    >{\bfseries\raggedright\arraybackslash}p{0.16\textwidth}
    >{\raggedright\arraybackslash}X
    >{\raggedright\arraybackslash}X
    >{\raggedright\arraybackslash}X}
\toprule
Aspetto &
\bfseries GDACS &
\bfseries CDSE/STAC &
\bfseries Copernicus CEMS \\
\midrule
Domanda &
Quali eventi richiedono attenzione? &
Quali acquisizioni satellitari coprono area e periodo? &
Quali attivazioni e prodotti di emergenza sono disponibili? \\
\addlinespace
Unità &
Evento e allerta &
Item e collezione satellitare &
Attivazione, area di interesse e prodotto \\
\addlinespace
Uso &
Monitoraggio di eventi &
Heatmap e bounding box dei prodotti &
Prodotti Rapid Mapping e Risk/Recovery \\
\addlinespace
Integrazione &
Server MCP GDACS &
Client STAC nel backend &
Server MCP CEMS \\
\bottomrule
\end{tabularx}
\caption{Ruoli delle fonti del percorso operativo. Marine e Drought alimentano
inoltre il catalogo di Discovery descritto nella Sezione~\ref{sec:fonti-discovery}.}
\label{tab:fonti-geospaziali}
\end{table}

Le tre sorgenti rispondono a domande diverse e possono essere combinate. Per
una richiesta su un'alluvione, GDACS fornisce gli eventi e il livello di
allerta; CEMS può indicare un'attivazione di cartografia; STAC può mostrare
quali acquisizioni satellitari sono presenti nell'area e nel periodo. Una
ricerca web aggiunge infine contesto recente e fonti narrative.

\section{Implicazioni per la progettazione dell'agente}

I server MCP e il client STAC convertono feed XML, risposte STAC e record CEMS
in dataset con campi stabili. In questo modo i client possono essere provati
separatamente dal modello, i rami indipendenti possono procedere in parallelo e
il guasto di una sorgente lascia intatti i risultati già disponibili.

Il workflow conserva semantica e provenienza di ogni record anche quando GDACS,
CDSE/STAC e CEMS sono interrogati nello stesso turno. Allo stesso modo, il grafo
mantiene distinti dataset, variabile, collezione e layer: una risorsa scoperta
indica quali dati potrebbero sostenere uno studio, il cui risultato richiede poi
un'analisi dedicata. Su questa separazione si basa l'architettura del capitolo
successivo.
